\documentclass[12pt,a4paper,oneside]{article}
\usepackage[brazil]{babel}
\usepackage[utf8]{inputenc}
\usepackage{olymp}

\setcounter{problem}{0}

\begin{document}

\begin{problem}{Caixa de ovos}{caixa}{2}
 
Bob está cansado dos tamanhos usados nas caixas de ovos. Uma caixa com 1 dúzia de ovos é geralmente de tamanho $2 \times 6$. Por que não $3 \times 4$? E uma caixa de 30 ovos, por que sempre $5 \times 6$? Não seria legal também ter os tamanhos $3 \times 10$ e $2 \times 15$?

Faça um programa para ajudar Bob a encontrar os tamanhos possíveis para caixas com diversas quantidades de ovos.

%\Notes
%
%Observações, se tiver.

\InputFile

A entrada contém apenas um inteiro $N$, o número de ovos na caixa. Restrições: $1 \leq N \leq 10000$.


\OutputFile

Seu programa deve escrever todas as possíveis configurações de caixa para exatamente $N$ ovos. Cada configuração deve ser escrita em uma linha no formado ``\texttt{L x C}'', sendo \texttt{L} o número de linhas da caixa e \texttt{C} o número de colunas (\texttt{x} é a letra xis minúscula).

\Notes
As configurações devem ser escritas em ordem crescente de número de linhas e não devem ser escritas configurações espelhadas. Por exemplo, para 30 ovos escreva ``\texttt{5 x 6}'' e não ``\texttt{6 x 5}''.

\Example

\begin{example}
\exmp{
30
}{
1 x 30
2 x 15
3 x 10
5 x 6
}%
\end{example}

\begin{example}
\exmp{
5
}{
1 x 5
}%
\end{example}

\begin{example}
\exmp{
25
}{
1 x 25
5 x 5
}%
\end{example}

\begin{example}
\exmp{
1000
}{
1 x 1000
2 x 500
4 x 250
5 x 200
8 x 125
10 x 100
20 x 50
25 x 40
}%
\end{example}


%Comentários sobre o exemplo, se necessário.

\end{problem}

\end{document}
